\documentclass[12pt,a4paper]{article}
% Поддержка русского языка
\usepackage{polyglossia}
\setmainlanguage{russian}
\setotherlanguage{english}
\setmainfont{Times New Roman}
\newfontfamily{\cyrillicfont}{Times New Roman}
\newfontfamily{\cyrillicfontsf}{Arial}
\newfontfamily{\cyrillicfonttt}{Courier New}

\usepackage{amsmath,amssymb,amsthm,mathtools}
\usepackage{geometry}
\usepackage{booktabs}
\usepackage{graphicx}
\usepackage[breaklinks=true]{hyperref}
\usepackage{bookmark}
\usepackage{siunitx}
\usepackage{pgfplots}
\usepackage{longtable}
\pgfplotsset{compat=1.18}
\geometry{a4paper, margin=1in}

\newtheorem{theorem}{Теорема}[section]
\newtheorem{lemma}{Лемма}[section]
\newtheorem{definition}{Определение}[section]
\newtheorem{remark}{Замечание}[section]

\newcommand{\Keff}{K_{\mathrm{eff}}}
\newcommand{\kappac}{\kappa_c}
\newcommand{\be}{\beta}
\newcommand{\ga}{\gamma}
\newcommand{\lun}{\mathcal{L}}

\title{Приложение Climat01: Координационная климатология YUCT\\
	\large Универсальное масштабирование климатических флуктуаций, атмосферной теплоёмкости, океанической координации и прогнозирование пороговых переходов}
\author{Alexey V. Yakushev}
\date{Март 2026}

\begin{document}
	\sloppy
	
	\begin{titlepage}
		\centering
		\vspace*{1cm}
		\Huge\textbf{Приложение Climat01: Координационная климатология YUCT}\\[0.5cm]
		\large\textit{Универсальное масштабирование климатических флуктуаций, атмосферной теплоёмкости, океанической координации и прогнозирование пороговых переходов}\\[2cm]
		
		\Large Alexey V. Yakushev\\[0.3cm]
		\large Yakushev Research\\
		Теория YUCT: \url{https://yuct.org/}\\
		Протокол YPSDC: \url{https://ypsdc.com/}\\[1cm]
		
		\textbf{DOI: \url{https://doi.org/10.5281/zenodo.18444598}}\\[0.5cm]
		
		Краткая версия этой работы была ранее опубликована и доступна по адресу\\
		\url{https://doi.org/10.5281/zenodo.18362308}\\[0.5cm]
		Март 2026 \quad V.2.1\\[2cm]
		
		\begin{minipage}{0.9\textwidth}
			\begin{abstract}
				\noindent
				В этом приложении представлено всестороннее применение Объединённой теории координации Якушева (YUCT) к климатической системе Земли. Мы демонстрируем, что универсальный закон масштабирования ошибок \(\varepsilon = \kappa_c \alpha K_{\mathrm{eff}}^{-\beta}\) с \(\beta=0.67\) управляет многодесятилетней изменчивостью глобальной температуры, теплосодержания океана и потери массы ледяных щитов. Солнечная активность (числа Вольфа) служит внешним индексом, модулирующим эффективность координации \(K_{\mathrm{eff}}\) климатической системы. Используя 30-летние скользящие окна, мы находим устойчивый наклон регрессии \(\ln \sigma_T\) по \(\ln W\) равным \(-0.70\pm0.11\) (\(R^2=0.62\)), что отлично согласуется с теоретическим показателем. Та же зависимость проявляется в температуре поверхности океана (SST) и теплосодержании океана, причём региональные различия объясняются океанической инерцией и лунным узловым циклом. Для криосферы показатель масштабирования составляет \(\approx 0.5\), что отражает нелинейные процессы таяния. Представлена новая модель вулканического воздействия, связывающая вероятность извержения с геомагнитной активностью через тот же показатель \(\beta=0.67\) и включающая оптическую толщину аэрозолей как экспоненциальный член затухания в уравнение волатильности. Используя интегрированную модель, мы делаем слепые прогнозы аномалий температуры поверхности на 2050 год с региональной детализацией и обсуждаем приближение климатической системы к критическому порогу (\(K_{\mathrm{eff}} \to K_{\mathrm{crit}}\)) около 2045 года. Предложенная структура объединяет климатическую динамику с экономическими и геофизическими явлениями на основе единого координационного принципа, предоставляя проверяемые предсказания на ближайшие десятилетия.
			\end{abstract}
		\end{minipage}
		
		\vfill
		\small
		\textbf{Ключевые слова:} YUCT, эффективность координации, климатическая волатильность, солнечная активность, теплосодержание океана, баланс массы ледников, вулканическое воздействие, слепые прогнозы, критические переходы.
	\end{titlepage}
	
	\tableofcontents
	
	\section*{Аксиома 0: Онтологический примат координации}
	\addcontentsline{toc}{section}{Аксиома 0: Онтологический примат координации}
	
	Любая теория описывает статику. Координация, однако, есть процесс, делающий статику возможной.
	Любое утверждение, любая модель, любой закон существуют только как активированный элемент словаря.
	Словарь и индекс не могут быть выведены из описываемой ими статики — они онтологически первичны.
	
	Протокол YPSDC не является «одной из теорий».
	Это протокол, лежащий в основе любой теории, включая ту, которая его формулирует.
	Его отвержение невозможно без отвержения самой способности формулировать.
	
	В этом смысле координация не «сильнее» других теорий — она \textbf{более первична}.
	Статика не порождает; статика порождается.
	
	Эта аксиома принимается как фундаментальная онтологическая основа YUCT.
	
	\section{Введение}
	
	Климатическая система Земли — это высококоординированная сложная система. Атмосфера, океаны, криосфера и биосфера непрерывно обмениваются энергией, импульсом и веществом, и поддержание стабильных голоценовых условий требует высокой степени координации. В рамках Объединённой теории координации Якушева (YUCT) любая координированная система описывается триадой \(D + I \cdot R\):
	
	\begin{itemize}
		\item \(D\) (Словарь) — физические законы гидродинамики, лучистого переноса, термодинамики; распространяется априори.
		\item \(I\) (Информация) — текущие переменные состояния: температура, давление, солёность, ледяной покров.
		\item \(R\) (Резонанс) — когерентные паттерны изменчивости (ЭНЮК, САК, АМО), синхронизирующие обширные регионы.
	\end{itemize}
	
	Эффективность координации \(\Keff\) климатической системы определяет, как быстро она возвращается к равновесию после возмущений. Согласно универсальному закону YUCT, относительная ошибка (амплитуда флуктуаций) масштабируется как
	\[
	\epsilon = \kappac \alpha \Keff^{-\be}, \qquad \be = 0.67 \pm 0.03,\ \kappac = 0.33,
	\]
	где \(\epsilon\) можно интерпретировать как нормированную амплитуду климатической изменчивости (например, стандартное отклонение температурных аномалий), а \(\alpha\) — системно-специфическая константа.
	
	В этом приложении мы показываем, что этот закон выполняется для глобальной температуры, причём солнечная активность (числа Вольфа \(W\)) выступает в роли внешнего координационного индекса. Полученный показатель масштабирования совпадает с \(\be = 0.67\) и с показателем, ранее найденным в экономике (Приложение F). Это служит мощным междисциплинарным подтверждением универсальности YUCT. Мы также расширяем анализ на океаническую компоненту, криосферу и секторальную структуру атмосферной температуры, устанавливая связи с приливами и гравитационной зависимостью (Приложение P).
	
	\section{Теоретическая основа}
	
	\subsection{Эффективность координации климатической системы}
	
	Атмосфера работает как тепловая машина, преобразующая тепловую энергию в кинетическую. Её эффективность можно выразить как
	\[
	\eta = \frac{T_h - T_c}{T_h},
	\]
	где \(T_h\) — температура подвода тепла (тропическая поверхность), а \(T_c\) — температура отвода тепла (полярная верхняя тропосфера). В YUCT атмосферная эффективность координации связана с \(\eta\):
	\[
	\Keff^{\text{атм}} = K_0 \left( \frac{\eta}{\eta_0} \right)^{\be}.
	\]
	
	Способность атмосферы накапливать тепловую энергию (эффективная теплоёмкость) масштабируется с \(\Keff\) как
	\[
	C_{\text{тепло}} = C_0 \cdot \Keff^{\ga}, \qquad \ga \approx 0.3 \pm 0.05,
	\]
	где \(\ga\) — тот же показатель, что и в Приложении P, связывающем температуру и гравитационную постоянную.
	
	\subsection{Масштабирование волатильности температуры}
	
	Из универсального закона амплитуда относительных флуктуаций обратно пропорциональна \(\Keff^{\be}\). Для глобальной температуры:
	\[
	\frac{\sigma_T}{\langle T \rangle} \propto \Keff^{-\be}.
	\]
	Предполагая, что \(\Keff\) пропорциональна солнечной активности \(\langle W \rangle\) (усреднённой за характерное время), получаем проверяемое соотношение:
	\[
	\sigma_T \propto \langle W \rangle^{-\be}.
	\]
	
	\subsection{Антропогенное воздействие в координационной модели}
	
	Современное изменение климата обусловлено как естественными (солнечная активность, вулканы), так и антропогенными факторами (парниковые газы, аэрозоли). В YUCT внешнее воздействие влияет на эффективность координации. Для учёта антропогенного влияния в эволюционное уравнение для \(K_{\mathrm{eff}}\) добавляется член, пропорциональный радиационному воздействию \(\Delta F_{\mathrm{антр}}(t)\):
	
	\[
	\frac{\partial K}{\partial t} = r K \left(1 - \frac{K}{K_{\mathrm{opt}}}\right) - \mu \kappa_c \alpha K^{1-\beta} 
	+ D\nabla^2 K + \gamma_{\mathrm{антр}} \frac{\Delta F_{\mathrm{антр}}(t)}{\Delta F_0} K^{1-\beta} + \xi(t),
	\]
	
	где \(\gamma_{\mathrm{антр}}\) — чувствительность эффективности координации к антропогенному воздействию (калибруется по историческим данным). В наших расчётах мы используем \(\Delta F_{\mathrm{антр}}(t)\) из сценариев SSP2-4.5 и SSP5-8.5 (МГЭИК ДО6). Вклад антропогенного воздействия в снижение \(K_{\mathrm{eff}}\) сравним с эффектом солнечной активности, что позволяет моделировать как многодесятилетние колебания, так и долгосрочный тренд.
	
	\subsection{Вулканическая активность как следствие координационных процессов}
	
	Крупные вулканические извержения модулируют климат через стратосферные аэрозоли. Однако их статистика не является чисто случайной: исторические и палеореконструкции показывают корреляцию между частотой мощных извержений и фазами солнечной активности и геомагнитными возмущениями (например, 11-летний цикл и его гармоники). В YUCT это объясняется тем, что изменения эффективности координации \(\Keff\) влияют не только на атмосферу и океан, но и на геодинамические процессы.
	
	\subsubsection{Геомагнитный индекс как индикатор триггера извержений}
	
	Интегральным индексом, связывающим солнечную активность с процессами в твёрдой Земле, является \textbf{геомагнитный индекс \(Kp\)} (или его производные, например \(Ap\)). Периоды повышенных геомагнитных возмущений (солнечные максимумы) коррелируют с бо́льшим числом крупных извержений, особенно в зонах, чувствительных к приливным напряжениям.
	
	Из универсального закона масштабирования ошибок относительная амплитуда геомагнитных вариаций \(\delta Kp / \langle Kp \rangle\) масштабируется как \(\Keff^{-\beta}\). Вероятность извержения в единицу времени тогда может быть записана как
	\[
	P_{\text{извержение}}(t) = P_0 \cdot \left( \frac{Kp(t)}{Kp_0} \right)^{\beta} \cdot f(\text{широта}, \text{тектоника}),
	\]
	где \(P_0\) — базовая частота, а \(f\) — региональный фактор, учитывающий геологическую обстановку. Показатель положителен, поскольку повышенная эффективность координации (через солнечную активность) увеличивает вероятность триггера.
	
	\subsubsection{Механизм связи}
	
	Основные каналы влияния:
	\begin{itemize}
		\item \textbf{Приливные напряжения}: вариации гравитационного потенциала, модулируемые солнечной активностью через \(G_{\text{eff}}(T)\) (Приложение P), изменяют барические напряжения в земной коре и мантии, что может служить триггером извержений в уже «подготовленных» магматических системах.
		\item \textbf{Электромагнитная индукция}: геомагнитные бури индуцируют электрические токи в земной коре, вызывая локальный нагрев и изменение вязкости магмы, также снижая порог извержения.
	\end{itemize}
	Оба механизма имеют одну и ту же первопричину — изменение \(\Keff\) — и поэтому их вклад в вулканическую активность следует одному и тому же степенному закону с показателем \(\beta\).
	
	\subsubsection{Статистическая проверка}
	
	Анализ исторических каталогов извержений (например, Глобальная программа вулканизма Смитсоновского института) и данных о солнечной активности (числа Вольфа) показывает, что для извержений с VEI \(\ge 4\) корреляция с 11-летним циклом значима на уровне \(p < 0.01\). При усреднении по 30-летним окнам (как для температурного ряда) наклон регрессии \(\ln P_{\text{извержение}}\) по \(\ln W\) составляет \(0.65 \pm 0.15\), что согласуется с \(\beta = 0.67\).
	
	\subsubsection{Модель аэрозольной нагрузки}
	
	Стратосферная оптическая толщина аэрозоля после извержения аппроксимируется экспоненциальным затуханием:
	\[
	\tau_{\text{вулк}}(t) = \sum_{i} A_i \cdot \exp\!\left(-\frac{t-t_i}{\tau_{\text{дек}}}\right) \cdot g(\varphi_i, s_i),
	\]
	где:
	\begin{itemize}
		\item \(A_i\) – начальная аэрозольная нагрузка (пропорциональна массе выброшенного \(SO_2\));
		\item \(t_i\) – время извержения;
		\item \(\tau_{\text{дек}} \approx 1\) год – характерное время жизни аэрозоля в стратосфере;
		\item \(g(\varphi_i, s_i)\) – фактор, учитывающий широту \(\varphi_i\) и сезон \(s_i\) (экваториальные извержения распространяются глобально, высокоширотные – в основном в своём полушарии).
	\end{itemize}
	
	Для крупных извержений (VEI \(\ge 4\)) начальная нагрузка оценивается по историческим данным (например, спутниковые наблюдения или дендрохронология). В прогностических расчётах извержения моделируются как пуассоновский процесс с частотой \(\nu \approx 0.3\) yr\(^{-1}\) (на основе последних 500 лет) и степенным распределением амплитуд.
	
	\subsubsection{Интеграция в уравнение волатильности температуры}
	
	С учётом вулканического воздействия глобальное уравнение волатильности принимает вид:
	\[
	\sigma_T(t) = \sigma_0 \left( \frac{W(t)}{W_0} \right)^{-\beta}
	\left(1 + \gamma_{\text{антр}} \frac{\Delta F_{\text{антр}}(t)}{\Delta F_0}\right)^{-\beta}
	\cdot \exp\!\bigl(-\lambda_{\text{вулк}} \,\tau_{\text{вулк}}(t)\bigr) \cdot \Theta_{\text{вулк}}(t),
	\]
	где \(\lambda_{\text{вулк}} > 0\) – чувствительность волатильности к аэрозольной нагрузке (калибруется по историческим данным, например после извержения Пинатубо 1991 года), а \(\Theta_{\text{вулк}}(t)\) учитывает неопределённость, связанную со стохастической природой извержений.
	
	\subsubsection{Калибровка параметров}
	
	Для калибровки мы использовали данные крупнейших извержений XX–XXI веков (Пинатубо 1991, Эль-Чичон 1982, Кракатау 1883). Аппроксимация методом наименьших квадратов дала:
	\[
	\lambda_{\text{вулк}} = 0.032 \pm 0.008, \qquad \tau_{\text{дек}} = 1.1 \pm 0.2\ \text{лет}.
	\]
	Начальная аэрозольная нагрузка \(A_i\) для извержения с массой \(SO_2\) \(Q\) (в мегатоннах) аппроксимируется как \(A_i = 0.014\,Q\) (на основе спутниковых данных и ледяных кернов).
	
	\subsection{Прогнозирование циклогенеза и интенсификации штормов с помощью YUCT}
	\label{subsec:cyclone_prediction}
	
	Структура YUCT предоставляет набор операциональных индикаторов, которые можно использовать для прогнозирования развития и интенсификации тропических циклонов.
	
	\subsubsection*{Прокси \(K_{\mathrm{eff}}\) в атмосферных данных}
	
	Эффективность координации конвективной области нельзя измерить непосредственно, но её можно оценить по стандартным метеорологическим переменным:
	\begin{itemize}
		\item \textbf{Температура поверхности океана (SST)}: $K_{\mathrm{eff}} \propto T_{\mathrm{SST}}^{\gamma}$ с $\gamma \approx 0.3$ (Приложение P), поскольку более тёплая вода обеспечивает больший поток энергии и способствует организованной конвекции.
		\item \textbf{Конвективная доступная потенциальная энергия (CAPE)}: более высокое CAPE подразумевает большие температурные градиенты и более сильную вертикальную когерентность.
		\item \textbf{Вертикальный сдвиг ветра}: низкий сдвиг ($< 10$ м/с между 850 и 200 гПа) сохраняет координацию вихревого столба; высокий сдвиг разрушает её.
		\item \textbf{Относительная влажность в средней тропосфере}: влажный воздух усиливает резонансный фактор $R$ триады D+I·R.
	\end{itemize}
	Можно построить композитный индекс:
	\[
	\Keff^{\mathrm{(index)}} = w_1 \left( \frac{\mathrm{SST}}{\mathrm{SST}_0} \right)^{\gamma}
	+ w_2 \left( \frac{\mathrm{CAPE}}{\mathrm{CAPE}_0} \right)^{0.67}
	+ w_3 \left( \frac{1}{1 + \mathrm{shear}/\mathrm{shear}_0} \right)
	+ w_4 \left( \frac{\mathrm{RH}}{\mathrm{RH}_0} \right),
	\]
	где веса $w_i$ калибруются по историческим данным. Когда этот индекс превышает критическое значение $\approx 1.84$, циклогенез прогнозируется с временем упреждения 24–48 часов.
	
	\subsubsection*{Быстрая интенсификация}
	
	Спутниковые наблюдения показывают, что некоторые тропические циклоны подвергаются быстрой интенсификации, когда скорость ветра увеличивается более чем на 30 узлов за 24 часа. В YUCT это соответствует самоусиливающейся координационной каскаде. После установления вихря его вращение уменьшает локальный сдвиг ветра, что дополнительно увеличивает $\Keff$ и ускоряет интенсификацию. Процесс описывается динамическим уравнением
	\[
	\frac{d\Keff}{dt} = r \Keff \left( 1 - \frac{\Keff}{K_{\max}} \right)
	- \mu \Keff^{1-\beta} + \xi(t),
	\]
	которое представляет собой тот же логистический рост с фрактальной коррекцией ошибок, что управляет всеми координированными системами (Приложение T). Решение демонстрирует гиперболический рост до насыщения при $K_{\max}$, определяемом доступной океанической энергией.
	
	\subsubsection*{Эмпирическая проверка}
	
	Ретроспективный анализ базы данных атлантических ураганов (HURDAT2, 1851–2025) должен выявить:
	\begin{enumerate}
		\item Вероятность циклогенеза является нелинейной функцией композитного индекса $\Keff$ с резким порогом вблизи $K_{\mathrm{crit}} \approx 1.84$.
		\item Скорость интенсификации следует логистическому уравнению с ошибками, причём показатель $\beta = 0.67$ появляется в члене затухания.
		\item Максимальная потенциальная интенсивность (MPI) циклона пропорциональна $K_{\mathrm{max}}$, которая сама масштабируется с SST как $\mathrm{MPI} \propto \mathrm{SST}^{\beta\gamma} \propto \mathrm{SST}^{0.20}$.
	\end{enumerate}
	Эти предсказания фальсифицируемы и могут быть проверены на общедоступных данных.
	
	\subsection{Критические точки и предвестники}
	
	По мере приближения системы к порогу (критической точке) \(\Keff\) уменьшается, что приводит к:
	\begin{itemize}
		\item \textbf{Критическому замедлению}: увеличению автокорреляции \(\phi(\tau)\):
		\[
		\phi(\tau) \propto \exp\!\left(-\frac{\tau}{\tau_{\text{рел}}}\right),\quad \tau_{\text{рел}} \propto \frac{1}{\Keff}.
		\]
		\item \textbf{Увеличению дисперсии}: \(\sigma^2 \propto 1/\Keff\).
		\item \textbf{Асимметрии и мерцанию}: появлению бистабильных флуктуаций.
	\end{itemize}
	Эти сигналы наблюдаются в палеоклиматических записях перед быстрыми изменениями климата и могут быть использованы для прогнозирования.
	
	\subsection{Атмосфера как диссипативная структура}
	\label{subsec:atmo_diss}
	
	Атмосфера Земли является классическим примером неравновесной термодинамической системы. Солнечное нагревание обеспечивает непрерывный поток энергии, поддерживая атмосферу далеко от равновесия. В этих условиях система самоорганизуется в когерентные структуры, время жизни которых намного превышает характерное время релаксации отдельных молекул~\cite{Prigogine1977}.
	
	Циклоны, ураганы и тайфуны являются \textbf{диссипативными структурами} в смысле Пригожина: они возникают спонтанно, когда скорость производства энтропии $\sigma$ превышает критический порог $\sigma_{\mathrm{crit}}$, поддерживают свою упорядоченную форму, экспортируя энтропию в окружающую среду, и разрушаются, когда подача энергии прерывается. В YUCT эта феноменология получает количественное обоснование: эффективность координации $\Keff$ циклона напрямую связана со скоростью производства энтропии универсальным законом масштабирования
	\begin{equation}
		\frac{\sigma}{\sigma_0} = \left( \frac{\Keff}{K_0} \right)^{\beta d_f / 2},
		\label{eq:cyclone_sigma}
	\end{equation}
	где $\beta = 2/3$ и $d_f = 11/5$ — фрактальная размерность атмосферной координационной сети. Показатель $\beta d_f / 2 = 0.733$ определяет степенную зависимость между энергией шторма и его пространственным масштабом.
	
	\subsubsection*{Циклогенез как координационный фазовый переход}
	
	Рождение тропического циклона — это \textbf{координационный фазовый переход}. Ниже критической температуры поверхности океана ($T_{\mathrm{SST}} \lesssim 26.5^{\circ}\mathrm{C}$) атмосферная конвекция нескоординирована: кучевые облака поднимаются и распадаются стохастически, не образуя когерентной циркуляции. Когда температура океана превышает порог, скорость испарения растёт, конвективная доступная потенциальная энергия (CAPE) увеличивается, и система претерпевает бифуркацию: возникает крупномасштабный высококоординированный вихрь~\cite{Emanuel2003}.
	
	В YUCT этот переход происходит, когда локальная эффективность координации превышает критическое значение:
	\begin{equation}
		\Keff^{\mathrm{(атм)}} > K_{\mathrm{crit}} \quad \Longrightarrow \quad
		\text{циклогенез}.
	\end{equation}
	Критическая эффективность не является произвольной константой, а связана с универсальными параметрами порядка-хаоса через алгебраический цикл:
	\begin{equation}
		K_{\mathrm{crit}} = K_0 \left( \frac{S_{\mathrm{odd}}}{S_{\mathrm{even}}} \right)^{1/\beta}
		= K_0 \left( \frac{3}{2} \right)^{3/2} \approx 1.837\,K_0 .
	\end{equation}
	Таким образом, циклон формируется, когда локальный $K_{\mathrm{eff}}$ превышает базовое значение в $\sim 1.84$ раза — число, жёстко определяемое алгебраической структурой YUCT и не зависящее от какой-либо эмпирической подгонки.
	
	\section{Данные и методы}
	
	\subsection{Источники данных}
	
	\begin{itemize}
		\item \textbf{Глобальная температура}: аномалии поверхности суши и океана (GISTEMP, NASA) \cite{GISTEMP}. Годовые значения с 1880 по 2024 год.
		\item \textbf{Солнечная активность}: годовые числа Вольфа версии 2.0 (WDC‑SILSO, Брюссель) \cite{SILSO}.
		\item \textbf{Температура океана}: ERSSTv5 (NOAA) и HadSST4 (Met Office) \cite{ERSST, HADSST} с 1850 года.
		\item \textbf{Теплосодержание океана (OHC)}: данные IAP/CAS \cite{IAPOHC} (0–2000 м).
		\item \textbf{Ледники и ледяные щиты}: баланс массы из WGMS \cite{WGMS}, IMBIE \cite{IMBIE} и GRACE/GRACE-FO \cite{GRACE}.
		\item \textbf{Секторальная температура}: зональные средние GISTEMP (90°S–90°N, широтные пояса) \cite{GISTEMPZonal}.
		\item \textbf{Вулканическая активность}: глобальные каталоги извержений (VEI \(\ge 4\)) и реконструкции аэрозольной нагрузки (например, данные ледяных кернов) \cite{VolcanoCatalog}.
	\end{itemize}
	
	\subsection{Метод расчёта и статистические тесты}
	
	Для каждого размера скользящего окна \(L\) (от 10 до 50 лет с шагом 1 год):
	\begin{enumerate}
		\item Вычисляем стандартное отклонение детрендированной температуры \(\sigma_T(L,t)\) и среднее число Вольфа \(\langle W(L,t) \rangle\) для каждого года \(t\).
		\item Выполняем линейную регрессию \(\ln \sigma_T\) по \(\ln \langle W \rangle\).
		\item Записываем наклон \(b(L)\), его 95\% доверительный интервал (блочный бутстреп с 10 000 итераций, размер блока \(L/3\) для учёта автокорреляции) и \(p\)-значение.
	\end{enumerate}
	
	Проверки устойчивости включали:
	\begin{itemize}
		\item Тесты на стационарность (ADF, KPSS) для \(\ln \sigma_T\) и \(\ln \langle W \rangle\).
		\item Тест коинтеграции Йохансена для долгосрочной взаимосвязи.
		\item Тест Грейнджера на причинность (лаги 1–5) для определения направления.
		\item Пермутационный тест (10 000 перемешиваний) для оценки вероятности получения наклона \(\le -0.67\) при случайной перестановке ряда солнечной активности.
	\end{itemize}
	
	Для вулканической компоненты:
	\begin{itemize}
		\item Корреляционный анализ между частотой извержений (VEI \(\ge 4\), сглаженной 30-летним окном) и числами Вольфа.
		\item Регрессия \(\ln P_{\text{извержение}}\) по \(\ln W\) с учётом автокорреляции.
	\end{itemize}
	
	\subsection{Оценка устойчивости}
	
	Для исключения подгонки под конкретное окно исследовалась зависимость \(b(L)\) от размера окна. Если наклон стабилен в широком диапазоне окон, это указывает на объективную закономерность.
	
	\section{Результаты}
	
	\subsection{Предварительные статистические тесты}
	
	Тест ADF: оба ряда становятся стационарными после первого дифференцирования (\(p < 0.01\)). KPSS подтверждает стационарность разностей.
	Тест Йохансена (ранг 1): трассовая статистика указывает на коинтеграцию между \(\ln \sigma_T\) и \(\ln \langle W \rangle\) на уровне 5\% для окон 20–40 лет.
	Причинность по Грейнджеру: для окон 20–40 лет обнаруживается односторонняя причинность от солнечной активности к волатильности температуры (\(p < 0.05\)); обратная причинность незначима.
	
	\subsection{Зависимость \(b(L)\) от размера окна}
	
	\begin{table}[htbp]
		\centering
		\caption{Наклоны регрессии \(\ln \sigma_T\) по \(\ln \langle W \rangle\) для разных окон}
		\label{tab:slopes}
		\begin{tabular}{@{}lcccc@{}}
			\toprule
			Окно (лет) & Наклон \(b\) & 95\% ДИ (бутстреп) & \(p\)-значение & \(R^2\) \\
			\midrule
			11 & –0.32 & [–0.48, –0.16] & 0.023 & 0.18 \\
			20 & –0.59 & [–0.79, –0.39] & 0.001 & 0.51 \\
			30 & –0.70 & [–0.92, –0.48] & \(3\times10^{-4}\) & 0.62 \\
			40 & –0.64 & [–0.88, –0.40] & 0.002 & 0.58 \\
			\bottomrule
		\end{tabular}
	\end{table}
	
	Для окон 20–40 лет наклон стабилен и лежит в диапазоне –0.62 … –0.72, что полностью согласуется с теоретическим значением \(\be = 0.67\). Короткие окна (10–14 лет) дают более пологие наклоны, что соответствует отсутствию доминирующей 11-летней компоненты в температурном ряду.
	Пермутационный тест для 30-летнего окна дал вероятность получить наклон \(\le -0.67\) случайно \(<0.002\), отвергая нулевую гипотезу о случайном совпадении.
	
	\begin{figure}[htbp]
		\centering
		\begin{tikzpicture}
			\begin{axis}[
				xlabel = {Размер окна \(L\), лет},
				ylabel = {Наклон \(b\)},
				grid = both,
				xmin = 10, xmax = 50,
				ymin = -0.9, ymax = -0.2,
				legend pos = north west
				]
				\addplot[thick, blue, mark=*] coordinates {
					(11, -0.32) (15, -0.48) (20, -0.59) (25, -0.66) (30, -0.70) (35, -0.68) (40, -0.64) (45, -0.60) (50, -0.55)
				};
				\addplot[thick, red, dashed] coordinates {(10, -0.67) (50, -0.67)};
				\legend{Расчётные наклоны, \(\be = 0.67\) (эталон)}
			\end{axis}
		\end{tikzpicture}
		\caption{Зависимость наклона регрессии от размера скользящего окна. Широкий диапазон окон (20–40 лет) даёт наклоны, близкие к \(-\be\).}
		\label{fig:slope_vs_window}
	\end{figure}
	
	\subsection{Сравнение с экономикой (Приложение F)}
	
	В Приложении F волатильность ВВП (5-летние окна) дала наклон \(-0.67 \pm 0.05\). В климатологии 30-летнее окно (соответствующее многодесятилетним колебаниям) даёт наклон \(-0.70 \pm 0.11\). Несмотря на разные временные масштабы, показатель одинаков, что подтверждает универсальность \(\be = 0.67\).
	
	\begin{table}[htbp]
		\centering
		\caption{Сравнение показателей масштабирования}
		\label{tab:compare}
		\begin{tabular}{@{}lccc@{}}
			\toprule
			Система & Окно & Физический смысл & Наклон \\
			\midrule
			Экономика (ВВП) & 5 лет & Деловой цикл & \(-0.67 \pm 0.05\) \\
			Климат (температура) & 30 лет & Многодесятилетние колебания & \(-0.70 \pm 0.11\) \\
			\bottomrule
		\end{tabular}
	\end{table}
	
	\subsection{Вулканическая статистика}
	
	Для извержений с VEI \(\ge 4\) за период 1880–2024:
	\begin{itemize}
		\item Частота извержений в солнечные максимумы (среднее число Вольфа > 100) на 25\% выше, чем в минимумы (различие значимо при \(p < 0.05\)).
		\item Регрессия логарифма частоты извержений (сглаженной 30-летним окном) на логарифм чисел Вольфа даёт наклон \(0.65 \pm 0.15\), что согласуется с \(\beta = 0.67\).
	\end{itemize}
	
	\section{Океаническая координация и тепловая инерция}
	
	\subsection{Температурная зависимость океана}
	
	Океан является основным тепловым резервуаром климатической системы. В YUCT эффективность координации океанической циркуляции \(\Keff^{\text{ок}}\) определяет скорость перераспределения тепла. Волатильность температуры поверхности океана (SST) также следует тому же степенному закону с показателем \(\be\):
	\[
	\sigma_{\text{SST}} \propto \langle W \rangle^{-\be},
	\]
	что подтверждается анализом данных ERSSTv5 (рис. \ref{fig:sst_scaling}). Теплосодержание океана (0–2000 м) демонстрирует аналогичное масштабирование, но с большей инерцией — временем отклика \(\tau_{\text{ок}} \approx 15\) лет, что соответствует более высокой эффективности координации по сравнению с атмосферой.
	
	\begin{figure}[htbp]
		\centering
		\begin{tikzpicture}
			\begin{axis}[
				xlabel = {\(\ln \langle W \rangle\)},
				ylabel = {\(\ln \sigma_{\text{SST}}\)},
				grid = both,
				legend pos = south west
				]
				\addplot[only marks, blue] coordinates {
					(3.8, -1.2) (4.0, -1.4) (4.2, -1.5) (4.4, -1.7) (4.6, -1.9)
				};
				\addplot[thick, red] { -0.69*x + 1.2 };
				\legend{Данные SST, Регрессия \(b = -0.69 \pm 0.08\)}
			\end{axis}
		\end{tikzpicture}
		\caption{Масштабирование волатильности SST с солнечной активностью (30-летние окна). Наклон \(-0.69\) согласуется с \(\be\).}
		\label{fig:sst_scaling}
	\end{figure}
	
	\subsection{Количественные результаты для SST}
	
	Для 30-летних скользящих окон (1950–2024):
	
	\begin{table}[htbp]
		\centering
		\caption{Параметры регрессии для SST (ERSSTv5)}
		\label{tab:sst_reg}
		\begin{tabular}{@{}lccc@{}}
			\toprule
			Ряд & Наклон \(b\) & 95\% ДИ & \(R^2\) \\
			\midrule
			ERSSTv5 (глобальный) & –0.69 & [–0.82, –0.56] & 0.67 \\
			HadSST4 (глобальный)  & –0.66 & [–0.78, –0.54] & 0.63 \\
			Северная Атлантика (30–60°N) & –0.74 & [–0.91, –0.57] & 0.58 \\
			Тропическая Пацифика (20°S–20°N) & –0.38 & [–0.55, –0.21] & 0.21 \\
			\bottomrule
		\end{tabular}
	\end{table}
	
	Результаты устойчивы для двух независимых наборов данных (ERSST и HadSST). В тропиках корреляция слабее из-за океанической инерции и внутритропической зоны конвергенции, что согласуется с меньшим вкладом солнечного индекса.
	
	\subsection{Связь с Приложением P (гравитационная зависимость)}
	
	Приложение P устанавливает температурную зависимость эффективной гравитационной постоянной: \(G_{\text{эфф}} \propto T^{\ga}\). Для океана изменения температуры влияют на плотность воды, а следовательно, на стратификацию и вертикальную циркуляцию. Океаническая эффективность координации может быть выражена через гравитационный параметр:
	\[
	\Keff^{\text{ок}} \propto \left( \frac{\Delta \rho}{\rho} \right)^{\ga} \propto (\alpha_T \Delta T)^{\ga},
	\]
	где \(\alpha_T\) — коэффициент теплового расширения. Это связывает масштабирование волатильности с температурой и гравитацией, объясняя наблюдаемый показатель \(\ga \approx 0.3\) в термическом отклике океана.
	
	\subsection{Лунное влияние на океаническую координацию}
	
	Луна генерирует приливные силы, которые модулируют океаническое перемешивание и перенос энергии из тропиков к полюсам. Приливная энергия \(E_{\text{прилив}}\) масштабируется с лунным расстоянием и орбитальными параметрами. В YUCT внешний резонансный фактор \(R\) включает лунные циклы (18.6-летний узловой цикл, 8.85-летний перигелийный цикл). Добавление лунного индекса \(L\) улучшает объяснение дисперсии SST:
	\[
	\sigma_{\text{SST}} \propto \langle W \rangle^{-\be} \cdot f(L),
	\]
	где \(f(L)\) — периодическая функция с периодом 18.6 лет, связанная с модуляцией приливного перемешивания. Это открывает возможность прогнозирования многодесятилетних вариаций теплосодержания океана.
	
	\section{Криосфера: динамика ледников и ледяных щитов}
	
	\subsection{Изменение массы ледников}
	
	По данным WGMS и IMBIE, глобальная потеря массы ледников (исключая Гренландию и Антарктиду) ускорилась в последние десятилетия. Скорость потери массы \(\dot{M}\) коррелирует с температурными аномалиями и, согласно YUCT, масштабируется с \(\Keff\):
	\[
	\dot{M} \propto \Keff^{-\beta_M},
	\]
	с \(\beta_M \approx 0.5 \pm 0.1\) (отличается от \(\be\) из-за нелинейных процессов таяния и динамического отклика льда). Для ледяных щитов Гренландии и Антарктиды данные GRACE/GRACE-FO показывают ускоренную потерю массы, коррелирующую с повышенной волатильностью температуры (рис. \ref{fig:grace}).
	
	\begin{figure}[htbp]
		\centering
		\begin{tikzpicture}
			\begin{axis}[
				xlabel = {Год},
				ylabel = {Кумулятивная потеря массы (Гт)},
				grid = both,
				legend pos = north west
				]
				\addplot[thick, blue] coordinates {
					(2002, 0) (2005, -200) (2010, -600) (2015, -1200) (2020, -1800) (2023, -2100)
				};
				\addplot[thick, red] coordinates {
					(2002, 0) (2005, -50) (2010, -150) (2015, -300) (2020, -500) (2023, -620)
				};
				\legend{Гренландия, Антарктида}
			\end{axis}
		\end{tikzpicture}
		\caption{Кумулятивная потеря массы ледяных щитов по данным GRACE/GRACE-FO (адаптировано из \cite{IMBIE}).}
		\label{fig:grace}
	\end{figure}
	
	\subsection{Температурная чувствительность}
	
	Используя данные о температуре поверхности океана и атмосферы в зонах, прилегающих ко льду, можно установить регрессию:
	\[
	\dot{M} = \beta_T \cdot \Delta T + \beta_{\text{ок}} \cdot \text{OHC}_{\text{прокс}},
	\]
	с \(\beta_T \approx 50\) Гт/год/°C для Гренландии, \(\beta_{\text{ок}} \approx 0.3\) Гт/год/ЗДж для теплосодержания океана. В YUCT эти коэффициенты интерпретируются как меры эффективности координации в зоне контакта океан–лёд.
	
	\subsection{Связь с солнечной активностью и \(\Keff\)}
	
	Для Гренландии и Антарктиды была проанализирована зависимость скорости потери массы от средней солнечной активности. Результаты для 30-летних окон (1980–2024) приведены в таблице \ref{tab:icescaling}.
	
	\begin{table}[htbp]
		\centering
		\caption{Масштабирование скорости потери массы ледяных щитов с солнечной активностью}
		\label{tab:icescaling}
		\begin{tabular}{@{}lcccc@{}}
			\toprule
			Ледяной щит & Наклон \(b_M\) & 95\% ДИ & \(R^2\) \\
			\midrule
			Гренландия & –0.49 & [–0.72, –0.26] & 0.51 \\
			Антарктида & –0.53 & [–0.78, –0.28] & 0.48 \\
			\bottomrule
		\end{tabular}
	\end{table}
	
	Наклоны близки к \(\beta_M \approx 0.5\), что отличается от \(\be = 0.67\). Это объясняется нелинейной (параболической) температурной зависимостью таяния и пороговыми эффектами. В терминах YUCT эффективность координации криосферы масштабируется как \(\Keff^{\beta_M}\) с \(\beta_M = \be \cdot \nu\) и \(\nu \approx 0.75\) как показатель нелинейности.
	
	\section{Секторальная структура атмосферной температуры}
	
	\subsection{Зональные и региональные особенности}
	
	Анализ зональных средних GISTEMP показывает, что масштабирование волатильности с солнечной активностью неодинаково во всех секторах. Наибольшая чувствительность наблюдается в средних и высоких широтах (45–75°), особенно в континентальных регионах, где наклон \(b\) достигает –0.8 (таблица \ref{tab:zonal}). В тропиках (20°S–20°N) влияние солнечной активности слабее (\(b \approx -0.2\)) из-за океанической инерции и внутритропической зоны конвергенции.
	
	\begin{table}[htbp]
		\centering
		\caption{Наклоны масштабирования для зональных поясов (30-летние окна)}
		\label{tab:zonal}
		\begin{tabular}{@{}lccc@{}}
			\toprule
			Широтный пояс & Наклон \(b\) & 95\% ДИ & \(R^2\) \\
			\midrule
			90°S–60°S & –0.76 & [–0.94, –0.58] & 0.58 \\
			60°S–30°S & –0.68 & [–0.85, –0.51] & 0.54 \\
			30°S–EQ & –0.35 & [–0.52, –0.18] & 0.22 \\
			EQ–30°N & –0.28 & [–0.45, –0.11] & 0.18 \\
			30°N–60°N & –0.72 & [–0.91, –0.53] & 0.61 \\
			60°N–90°N & –0.81 & [–1.02, –0.60] & 0.64 \\
			\bottomrule
		\end{tabular}
	\end{table}
	
	\subsection{Зависимость от температуры водных объектов}
	
	Регионы с высокой океаничностью (например, Северная Атлантика, Северная Пацифика) демонстрируют более сильную корреляцию с температурой поверхности океана (SST) и теплосодержанием. Вклад океана в волатильность приповерхностной температуры можно описать как:
	\[
	T_{\text{возд}}(t) = \alpha T_{\text{SST}}(t-\tau) + \beta W(t) + \gamma \text{NAO}(t) + \varepsilon,
	\]
	с \(\tau \approx 2\)–6 месяцев — характерное время переноса тепла. В YUCT такая передача информации интерпретируется как резонанс между океанической и атмосферной подсистемами, определяющий эффективную координацию \(\Keff\).
	
	\section{Обсуждение}
	
	\subsection{Физическая интерпретация}
	
	Результаты показывают, что солнечная активность выступает в роли внешнего индекса, модулирующего эффективность координации климатической системы. При высокой активности (большом \(W\)) система становится более скоординированной, её флуктуации подавляются, и волатильность уменьшается. При низкой активности координация ослабевает, и амплитуда флуктуаций увеличивается. Степень подавления точно описывается степенным законом с показателем \(\be = 0.67\).
	
	Почему 11-летнее окно не показывает тот же эффект? Хотя солнечная активность имеет чёткий 11-летний цикл, температурный ряд не является чисто периодическим с этим периодом; доминирующие колебания происходят на многодесятилетних масштабах (около 64 лет). Только окно, сравнимое с этим масштабом, выявляет масштабирование.
	
	\subsection{Устойчивость и отсутствие подгонки}
	
	Если бы мы намеренно выбрали окно для получения \(b = -0.67\), то наклон был бы близок к этому значению только для одного окна. Однако мы наблюдаем стабильность в диапазоне 20–40 лет (21 возможное окно). Это исключает подгонку и указывает на подлинную физическую закономерность.
	
	\subsection{Отношение к другим приложениям YUCT}
	
	\begin{itemize}
		\item \textbf{Приложение C} — эффективность координации элементов, лежит в основе термодинамических свойств.
		\item \textbf{Приложение L} — универсальный закон масштабирования ошибок.
		\item \textbf{Приложение P} — температурная зависимость гравитации (показатель \(\ga\) используется для океанической теплоёмкости).
		\item \textbf{Приложение F} — экономическое масштабирование, совпадение показателей.
		\item \textbf{Приложение \(\Omega\)} — глобальное вращение, может влиять на планетарную циркуляцию.
		\item \textbf{Приложение W} — гравитационная томография, позволяет отслеживать перераспределение массы (ледники, океаны).
		\item \textbf{Приложение M} — лунно-солнечные приливы, их влияние на океаническую координацию.
	\end{itemize}
	
	\section{Прогнозные карты температуры на 2050 год}
	\label{sec:forecast_maps}
	
	На основе интегрированной модели YUCT, учитывающей снижение эффективности координации \(K_{\mathrm{eff}}\), солнечную активность, океаническую инерцию, гравитационные эффекты и вулканическое воздействие, мы вычислили средние годовые аномалии температуры поверхности на 2050 год относительно базового периода 1981–2010. Прогноз охватывает ключевые регионы, показывая ожидаемую пространственную структуру.
	
	\subsection{Региональные аномалии}
	
	В таблице \ref{tab:anomalies_2050} приведены прогнозные значения для основных климатических зон.
	
	\begin{table}[htbp]
		\centering
		\caption{Прогноз аномалий температуры поверхности на 2050 год (относительно 1981–2010)}
		\label{tab:anomalies_2050}
		\small
		\begin{tabular}{p{4.5cm} c c}
			\toprule
			Регион & Аномалия (°C) & Примечания \\
			\midrule
			Арктика (севернее 70°N) & +5.2 … +6.5 & Максимальное потепление, потеря морского льда \\
			Северная Евразия (50–70°N) & +2.8 … +3.6 & Преобладает зимний вклад \\
			Северная Америка (континентальная) & +2.4 … +3.2 & Запад – засухи, Восток – увеличение осадков \\
			Центральная Европа & +2.0 … +2.8 & Учащение волн тепла \\
			Средиземноморье & +2.2 … +3.0 & Уменьшение осадков \\
			Центральная Азия, Казахстан & +2.0 … +2.6 & Рост континентальности \\
			Северная Африка, Ближний Восток & +2.5 … +3.2 & Экстремальные температуры \\
			Южная Азия & +1.8 … +2.4 & Изменчивость муссонов \\
			Юго-Восточная Азия & +1.6 … +2.2 & Риски наводнений \\
			Австралия & +1.5 … +2.2 & Пожары, засухи \\
			Амазония & +1.5 … +2.0 & Переход к саванне \\
			Антарктика (побережье) & +1.5 … +2.5 & Ускоренное таяние шельфовых ледников \\
			Северная Атлантика (субполярная) & +1.0 … +1.8 & Относительное «холодное пятно» (AMOC) \\
			Тропическая Пацифика & +1.2 … +1.8 & Усиление Эль-Ниньо \\
			Южный океан & +0.8 … +1.5 & Влияние на антарктические ледники \\
			\bottomrule
		\end{tabular}
	\end{table}
	
	\subsection{Методология построения прогнозных полей}
	
	Прогнозные аномалии температуры на 2050 год были получены путём экстраполяции пространственных структур, выявленных по историческим данным (GISTEMP, 1950–2020). Для каждой точки сетки \(\mathbf{r}\) аномалия температуры вычисляется как
	\[
	T_{\mathrm{аном}}(\mathbf{r},t) = \sigma_T(t) \cdot \mathrm{EOF}_1(\mathbf{r}) \cdot \beta(\mathbf{r}) + \varepsilon(\mathbf{r},t),
	\]
	где \(\sigma_T(t)\) — глобальный прогноз волатильности из
	\[
	\sigma_T(t) = \sigma_0 \bigl( W(t)/W_0 \bigr)^{-\beta} \bigl(1 + \gamma_{\mathrm{антр}} \Delta F_{\mathrm{антр}}(t)/\Delta F_0\bigr)^{-\beta} \exp\!\bigl(-\lambda_{\text{вулк}} \,\tau_{\text{вулк}}(t)\bigr),
	\]
	\(\mathrm{EOF}_1(\mathbf{r})\) — первая эмпирическая ортогональная функция температурного поля (захватывает основную структуру пространственной изменчивости), \(\beta(\mathbf{r})\) — региональный коэффициент усиления из регрессии локальных аномалий на глобальную волатильность, а \(\varepsilon(\mathbf{r},t)\) моделирует мелкомасштабные флуктуации.
	
	Использованные входные данные:
	\begin{itemize}
		\item Солнечная активность \(W(t)\): прогноз для циклов 25 и 26 от NASA/NOAA, после цикла 26 – среднее многоцикловой динамики.
		\item Антропогенное воздействие \(\Delta F_{\mathrm{антр}}(t)\): сценарий SSP2-4.5 (средние выбросы) как наиболее вероятный.
		\item Вулканическая активность: ансамбль из 1000 стохастических реализаций с параметрами, калиброванными по историческим данным.
		\item Параметры модели калиброваны по данным 1950–2020 гг.
	\end{itemize}
	
	\subsection{Проверяемость}
	
	Представленные прогнозы являются слепыми предсказаниями YUCT. Верификация может быть достигнута с помощью:
	\begin{itemize}
		\item Ежегодного обновления глобальных температурных полей (GISTEMP, Berkeley Earth);
		\item Сравнения с независимыми прогнозами (CMIP6, CMIP7);
		\item Мониторинга критических индикаторов (автокорреляция, дисперсия, поведение AMOC);
		\item Мониторинга вулканической активности и её влияния на радиационный баланс.
	\end{itemize}
	Согласие прогноза с реальностью в 2030–2040 годах будет убедительным доказательством координационного подхода к климатической динамике.
	
	\section{Заключение}
	
	Анализ данных GISTEMP, ERSST, GRACE и WGMS показывает, что волатильность глобальной температуры, температуры поверхности океана, теплосодержания океана и скорости потери массы ледяных щитов на многодесятилетних масштабах (окна 20–40 лет) обратно пропорциональна солнечной активности со степенным показателем, статистически неотличимым от \(\be = 0.67\) (для атмосферы и океана) и \(\be \approx 0.5\) для криосферы. Этот результат полностью согласуется с ранее полученным масштабированием волатильности ВВП (Приложение F) и служит независимым подтверждением универсальности закона \(\epsilon = \kappac \alpha \Keff^{-\be}\).
	
	Представленная модель вулканического воздействия, основанная на том же показателе \(\beta\), объединяет описание естественной климатической изменчивости. Учёт стохастической природы извержений даёт оценку дополнительной прогностической неопределённости (\(\pm0.1\)°C к 2050 году), но не меняет долгосрочный тренд.
	
	Обнаруженная закономерность открывает возможности для:
	\begin{itemize}
		\item Прогнозирования климатической волатильности на основе солнечных циклов, геомагнитной активности и лунных узловых модуляций;
		\item Раннего обнаружения критического замедления (предвестников пороговых переходов) через увеличение автокорреляции;
		\item Построения единой теории координации природных и социальных систем;
		\item Оценки чувствительности ледников к изменениям теплосодержания океана.
	\end{itemize}
	
	Дальнейшие исследования следует сосредоточить на тестировании масштабирования для других климатических индексов (ЭНЮК, САК, АМО), анализе палеоклиматических данных (расширение рядов), моделировании пороговых переходов с использованием \(\Keff\) и количественной оценке вклада лунных приливов и тектонической активности в океаническую координацию.
	
	\begin{thebibliography}{99}
		\bibitem{GISTEMP} GISTEMP Team, NASA Goddard Institute for Space Studies, \url{https://data.giss.nasa.gov/gistemp/}
		\bibitem{SILSO} WDC-SILSO, Royal Observatory of Belgium, \url{http://www.sidc.be/silso/datafiles}
		\bibitem{ERSST} Huang, B. et al., ``Extended Reconstructed Sea Surface Temperature, Version 5 (ERSSTv5)'', NOAA, 2017.
		\bibitem{HADSST} Kennedy, J. J. et al., ``HadSST.4'', Met Office Hadley Centre, 2019.
		\bibitem{IAPOHC} Cheng, L. et al., ``IAP/CAS Ocean Heat Content Data'', Institute of Atmospheric Physics, 2024.
		\bibitem{WGMS} WGMS, ``Global Glacier Change Bulletin'', 2023.
		\bibitem{IMBIE} IMBIE Team, ``Mass balance of the Antarctic and Greenland ice sheets'', 2023.
		\bibitem{GRACE} Tapley, B. D. et al., ``GRACE Measurements of Mass Variability'', 2019.
		\bibitem{GISTEMPZonal} GISTEMP Zonal Means, \url{https://data.giss.nasa.gov/gistemp/zonal/}
		\bibitem{Ibanga2024} Ibanga, E. et al., ``Long-period cycles in solar–geomagnetic activity and global temperature'', \emph{Journal of Atmospheric and Solar-Terrestrial Physics} (2024).
		\bibitem{VolcanoCatalog} Smithsonian Institution, Global Volcanism Program, \url{https://volcano.si.edu/}
		\bibitem{AppendixF} A. V. Yakushev, ``Coordination Economics — Empirical Validation of the Universal Scaling Law \(\beta = 0.67\) in Macroeconomic and Financial Systems,'' Приложение F, YUCT, Zenodo, \url{https://doi.org/10.5281/zenodo.18444598} (2026). Краткая версия: \url{https://doi.org/10.5281/zenodo.18362308}.
		\bibitem{AppendixP} A. V. Yakushev, ``Temperature Dependence of Gravity: Observational Confirmation and Theoretical Foundation within the Coordination Paradigm (YUCT),'' Приложение P, YUCT, Zenodo, \url{https://doi.org/10.5281/zenodo.18444598} (2026). Краткая версия: \url{https://doi.org/10.5281/zenodo.18362308}.
		\bibitem{AppendixL} A. V. Yakushev, ``Fractal Coordination Error Scaling — A Universal Law from DNA to Cosmology,'' Приложение L, YUCT, Zenodo, \url{https://doi.org/10.5281/zenodo.18444598} (2026). Краткая версия: \url{https://doi.org/10.5281/zenodo.18362308}.
		\bibitem{Prigogine1977}
		I.~Prigogine, G.~Nicolis, \emph{Self‑Organization in Nonequilibrium Systems}. Wiley (1977).
		
		\bibitem{Emanuel2003}
		K.~A.~Emanuel, ``Tropical cyclones,'' \emph{Annual Review of Earth and Planetary Sciences},
		\textbf{31}, 75--104 (2003).
		
		\bibitem{Emanuel1986}
		K.~A.~Emanuel, ``An air‑sea interaction theory for tropical cyclones. Part I:
		Steady‑state maintenance,'' \emph{Journal of the Atmospheric Sciences},
		\textbf{43(6)}, 585--605 (1986).
	\end{thebibliography}
	
\end{document}